\documentclass{article}
\usepackage[UTF8]{ctex}
\usepackage{amsmath,amssymb,amsfonts}
\usepackage{enumitem}
\usepackage{geometry}
\geometry{a4paper, margin=1in}
% === 错题本排版优化（不参与题目标识） ===
\setlength{\parindent}{0pt}
\setlength{\parskip}{0.35em plus 0.10em minus 0.08em}
\setlength{\emergencystretch}{3em}
\setlist{itemsep=0.25em, topsep=0.40em, parsep=0pt}
\allowdisplaybreaks[2]
% === 错题本排版优化结束 ===
\begin{document}
\section*{中国矿业大学(北京)2024年数学分析真题}
\section*{25．中国矿业大学（北京）2024年数学分析试卷}

\subsection{题目1}
1．（10 分）计算函数极限 \(\lim _{x \rightarrow 0} \frac{(1+x)^{\frac{2}{x}}-e^{2}(1-\ln (1+x))}{x^{2}}\) ．

\subsection{题目2}
2．（10 分）计算不定积分 \(\int \frac{x e^{x}}{\sqrt{e^{x}+1}} \mathrm{~d} x\) ．

\subsection{题目3}
3．（10 分）函数 \(z=z(x, y)\) 是由方程 \(\frac{x}{z}=\ln \frac{z}{y}\) 确定的隐函数，计算 \(\frac{\partial^{2} z}{\partial x \partial y}\) ．

\subsection{题目4}
4．（10 分）计算二重积分 \(\iint_{D} e^{-(x+y)^{2}} \mathrm{~d} x \mathrm{~d} y\) ，其中 \(D\) 是由 \(x+y=1, x=0, y=x\) 所围成的区域．

\subsection{题目5}
5．（10 分）计算曲线积分 \(\oint_{L} \frac{3 x \mathrm{~d} x-4 y \mathrm{~d} y}{3 x^{2}+4 y^{2}}\) ，其中 \(L\) 是 \(x^{2}+y^{2}=1\) 的正向边界．

\subsection{题目6}
6．（10 分）计算曲面积分

\[
\iint_{\sum} y^{3} \mathrm{~d} z \mathrm{~d} x+(y+z) \mathrm{d} x \mathrm{~d} y
\]

其中 \(\sum\) 为 \(z=x^{2}+y^{2}\) 被 \(z=0, z=1\) 所截部分的外侧。

\subsection{题目7}
7．（15 分）求函数 \(f(x)=\arctan \frac{1-x^{2}}{1+x^{2}}\) 的单调区间，极值，凹凸区间，拐点和渐近线．

\subsection{题目8}
8．（15 分）函数 \(f(x, y)=x^{2}+(y-1)^{2}(x \neq 0)\) 在 \(\frac{x^{2}}{a^{2}}+\frac{y^{2}}{b^{2}}=1\) 条件下有最小值 1 ，且椭圆的面积最小，求 \(a, b\) 的值。

\subsection{题目9}
9．（15 分）判断反常积分 \(\int_{0}^{+\infty} \frac{e^{\sin x} \sin 2 x}{x^{p}} \mathrm{~d} x(p>0)\) 的敛散性，并讨论条件收敛和绝对收敛．

\subsection{题目10}
10．（15 分）函数 \(f(x), g(x)\) 在区间 \([a, b]\) 上二阶可导，且 \(f(a)=g(a), f(b)=g(b), f(x), g(x)\) 在\((a, b)\) 内有相同的最大值．证明：存在 \(\xi \in(a, b)\) ，使得 \(f^{\prime \prime}(\xi)=g^{\prime \prime}(\xi)\) ．

\subsection{题目11}
11．（15 分）设 \(f(x)=\sum_{n=1}^{\infty} \frac{\cos n x}{\sqrt{n^{3}+n}}, g(x)\) 是 \(f(x)\) 的原函数，且 \(g(0)=0\) ，证明：

\begin{itemize}
  \item[（1）] \(f(x)\) 在 \(\mathbb{R}\) 上连续．
  \item[（2）] \(\frac{\sqrt{2}}{2}-\frac{1}{15}<g\left(\frac{\pi}{2}\right)<\frac{\sqrt{2}}{2}\) ．
\end{itemize}

\subsection{题目12}
12．（15 分）叙述实数系的闭区间套定理和确界原理，并用闭区间套定理证明确界原理．
\end{document}
